\chapter{探测器究竟读到了电子态的哪一部分}
\label{chap:feqo-measurement-inference}

上一章的时间密度包含边带之间的相位，而普通能谱仪只按能量分箱。若两个电子态有相同人口、不同相位，单次能谱会完全相同。为了说清探测器能读什么，先把一次测量写成从密度算符到结果概率的映射，再加入有限分辨率和计数涨落。

\section{从 Born 概率到有限分辨的能谱仪}
\label{chap:feqo-measurement}

对理想能量基 \(|\ell\rangle\)，投影 \(P_\ell=|\ell\rangle\langle\ell|\) 满足 \(P_\ell\geq0\)、\(\sum_\ell P_\ell=I\)。Born 规则给 \(p_\ell=\operatorname{Tr}(\rho P_\ell)=\rho_{\ell\ell}\)。因为 \(P_\ell\) 是对角的，任意非对角元都不会进入这次测量。即使上一章的时间波形有尖峰，理想能量人口仍只读 \(|c_\ell|^2\)。

有限分辨率把一个真实能量 \(E_\ell\) 分配到一系列读数 \(y\)。令 \(R(y|E_\ell)\geq0\) 是条件概率密度，满足 \(\int R(y|E_\ell)\dd y=1\)。把仪器响应放进测量算符，定义
\begin{equation}
M_y=\sum_\ell R(y|E_\ell)P_\ell,
\qquad p(y|\rho)=\operatorname{Tr}(\rho M_y)
=\sum_\ell R(y|E_\ell)\rho_{\ell\ell}.
\label{eq:feqo-energy-povm}
\end{equation}
每个 \(M_y\) 正，且 \(\int M_y\dd y=I\)。满足这两个条件的一族算符叫正算符值测量（positive operator-valued measure，POVM）；它只规定结果概率，不自动规定电子在读出以后怎样变化。响应 \(R\) 可以是高斯函数，例如 \(R(y|E)=(2\pi\sigma_D^2)^{-1/2}e^{-(y-E)^2/(2\sigma_D^2)}\)，其中 \(\sigma_D\) 是探测器能量分辨率，不是入射电子的能散。

这族算符的正性可直接检验，而无需把“有限分辨”当成一句直觉话。任取电子态 \(|\chi\rangle=\sum_\ell\chi_\ell|\ell\rangle\)，有 \(\langle\chi|M_y|\chi\rangle=\sum_\ell R(y|E_\ell)|\chi_\ell|^2\geq0\)；再对 \(y\) 积分，每个 \(R\) 的归一化给 \(\sum_\ell|\chi_\ell|^2\)。但 \(M_y\) 仍在能量基对角，所以把分辨率做得无限好也不能凭同一设置恢复相位；分辨率与相位可识别性是不同问题。

用相邻两条能量相差 \(\Delta E=\hbar\omega\) 的等先验边带，可以把“可分辨”算成一个数。若两条边带的读数都服从标准差 \(\sigma_D\) 的高斯响应，按读数是否超过两中心的中点分类，两侧的错误概率相同，都是
\[
P_{\rm err}=\frac12\operatorname{erfc}
\!\left(\frac{\Delta E}{2\sqrt2\,\sigma_D}\right),
\qquad
\operatorname{erfc}(x):=\frac2{\sqrt\pi}\int_x^\infty e^{-u^2}\dd u.
\]
推导只需把较低能量那一峰从中点积分到 \(+\infty\)，再作变量 \(u=(y-E_0)/(\sqrt2\sigma_D)\)。当 \(\sigma_D\ll\Delta E\) 时误判趋零；当 \(\sigma_D\gg\Delta E\) 时它趋向 \(1/2\)，测量对两通道近乎没有区分力。这个判错率与电子态本身是否纯、是否相干没有直接等号。

若第 \(\ell\) 通道的检出效率为 \(\eta_\ell\in[0,1]\)，可把 \(M_y\) 中的每项改成 \(\eta_\ell R(y|E_\ell)P_\ell\)，并增加一个“未检出”效应 \(M_\varnothing=\sum_\ell(1-\eta_\ell)P_\ell\)。这时 \(M_\varnothing+\int M_y\dd y=I\)，所以损失的计数没有被偷偷改写成不归一化概率。

把响应宽度和有限窗口放进同一张数值图，这两类偏差的差别就一目了然。\cref{fig:electron-detector-forward} 左图取同一个 \(|\beta|=1.20\) 的真实边带谱，只把高斯响应宽度由 \(\sigma_E=0.292\,\si{\electronvolt}\) 换成 \(0.800\,\si{\electronvolt}\)：两组的边带权重完全相同，差别只是相邻 \(\ell\) 的读数是否被抹平。右图把读数截在半宽 \(W\) 的窗口内，却仍用完整谱的方差公式反演，得到的 \(|\beta|\) 在 \(W/(\hbar\omega)=0.7\) 时只有 \(0.42\)，即使开到 \(2\) 也只回到 \(0.90\)，要一直开到 \(5\) 才收敛到虚线标出的真值 \(1.20\)。窗口截断造成的偏差是前向模型问题，不能记到计数噪声头上。

\begin{figure}[htbp]
\centering
\includegraphics[width=0.94\linewidth]{electron_detector_forward.pdf}
\caption{同一 \(|\beta|=1.20\) 电子谱经不同仪器分辨率与有限能量窗口。左图横轴为读数偏离 \(\varepsilon_d\)，纵轴为读数密度，两条曲线取响应宽度 \(\sigma_E=0.292\,\si{\electronvolt}\) 与 \(0.800\,\si{\electronvolt}\)，真实边带权重不变。右图横轴为半窗宽 \(W/(\hbar\omega)\)，纵轴为反演出的 \(|\beta|\)，虚线为真耦合 \(|\beta|=1.20\)，曲线是对条件数据直接套用完整方差公式的结果。}
\label{fig:electron-detector-forward}
\end{figure}

二维小矩阵可以同时检验正性和条件更新。取 \(M_+=\operatorname{diag}(0.8,0.2)\)、\(M_-=\operatorname{diag}(0.2,0.8)\)，则二者本征值非负且和为 \(I\)。一种具体测量装置可用 \(K_\pm=\sqrt{M_\pm}\) 描述结果后的态更新。对 \(|+\rangle=(|0\rangle+|1\rangle)/\sqrt2\)，\(p_+=\langle+|M_+|+\rangle=1/2\)；记录 \(+\) 后，归一化态为 \(\sqrt{0.8}|0\rangle+\sqrt{0.2}|1\rangle\)。同一 \(M_+\) 也可能由别的测量操作实现，故只知结果概率时不能擅自选择这一个条件态。

若要读相位，可在能量测量前施加已知酉变换 \(U_s\)。设置 \(s\) 下的概率变为
\begin{equation}
p(y|s,\rho)=\operatorname{Tr}\bigl(U_s\rho U_s^\dagger M_y\bigr).
\label{eq:feqo-phase-readout}
\end{equation}
变换把原来的非对角元旋转到新的能量人口。参考场相位、两次作用间的传播距离和收集角可提供不同设置；每一种设置必须先标定，才能将记录差异解释为电子相位。

\begin{exercise}
证明两个态 \((|0\rangle+|1\rangle)/\sqrt2\) 与 \((|0\rangle+\ii|1\rangle)/\sqrt2\) 的理想能量谱相同，并构造一个投影到 \((|0\rangle+e^{\ii\theta}|1\rangle)/\sqrt2\) 的设置区分它们。
\end{exercise}


知道了测量概率以后，才可以问一个反问题：从这些结果能否恢复未知场或未知态？答案由前向映射的可识别性决定；优化算法只能在已有信息中寻找参数。

\section{相位什么时候能由计数恢复}
\label{chap:feqo-inversion}

先给一个可手算的可识别性例子。两边带电子态写为
\[
\rho(r,\phi)=\frac12
\begin{pmatrix}1&r e^{-\ii\phi}\\r e^{\ii\phi}&1\end{pmatrix},
\qquad 0\leq r\leq1.
\]
两个本征值是 \((1\pm r)/2\)，故这确是正、迹为一的密度矩阵。能量投影对任意 \(r,\phi\) 都给 \(p_0=p_1=1/2\)；无论重复多少次，单一能谱都无法恢复相位。若增加投影到 \(|+_\theta\rangle=(|0\rangle+e^{\ii\theta}|1\rangle)/\sqrt2\) 的参考设置，则
\begin{equation}
p_+(\theta)=\langle+_\theta|\rho|+_\theta\rangle
=\frac{1+r\cos(\phi-\theta)}2.
\label{eq:feqo-reference-phase-probability}
\end{equation}
取 \(\theta=0,\pi/2\)，两次测量分别给 \(r\cos\phi\) 与 \(r\sin\phi\)，从而在 \(r>0\) 时恢复 \(r,\phi\)。当 \(r=0\) 时相位根本没有物理定义，任意测量都不能“恢复”它。这比只说“多设置有用”更明确：新增设置补上了前向映射缺少的独立方向。

两设置还允许直接计算实验难度。把两次“\(+\)”概率记为 \(p_0=(1+r\cos\phi)/2\)、\(p_{\pi/2}=(1+r\sin\phi)/2\)。它们对 \((r,\phi)\) 的导数矩阵为
\[
J=\frac12
\begin{pmatrix}
\cos\phi&-r\sin\phi\\
\sin\phi&r\cos\phi
\end{pmatrix},
\qquad \det J=\frac r4.
\]
所以 \(r>0\) 时，任意 \(\phi\) 的微小变化都留下可区分的一阶计数变化。为量化这种变化，先定义一次观测的\emph{得分}为对数概率的参数导数 \(S_a=\partial_a\log P_\vartheta(x)\)，定义 Fisher 信息矩阵为 \(I_{ab}=\mathbb E(S_aS_b)\)。若每个设置有 \(\nu_s\) 个独立入射电子，并同时记录“\(+\)”与“\(-\)”两类计数，那么单次“\(+\)”结果的得分是 \((\partial_a p_s)/p_s\)，单次“\(-\)”结果的是 \(-(\partial_a p_s)/(1-p_s)\)。按两种结果的概率加权，再对独立观测求和，得到这一二项实验的信息
\[
I_{ab}=\sum_{s\in\{0,\pi/2\}}
\frac{\nu_s\,(\partial_a p_s)(\partial_b p_s)}{p_s(1-p_s)},
\qquad 0<p_s<1.
\]
例如 \(r\ll1\) 且两次曝光同为 \(\nu\) 时，\(p_s(1-p_s)=1/4+O(r^2)\)，由两行导数的正交性可得 \(I_{rr}=\nu+O(\nu r^2)\)、\(I_{\phi\phi}=\nu r^2+O(\nu r^4)\)。估计相位的方差尺度因此随 \(r^{-2}\) 恶化。这里的“方差尺度”并非凭直觉猜测，而是来自下述可以直接证明的界。数学上 \(r>0\) 已可识别，实验上微弱相干仍可能需要很大曝光；这是“能否唯一求出”与“能否稳定求出”的区别。

\begin{proposition}[计数实验的 Cram\'er--Rao 界]
设参数 \(\vartheta\) 在开区间中，离散结果 \(x\) 的概率 \(P_\vartheta(x)>0\) 可微，且可把有限或绝对收敛的求和与求导交换。记得分 \(S_\vartheta(x)=\partial_\vartheta\log P_\vartheta(x)\) 和信息 \(I(\vartheta)=\mathbb E_\vartheta S_\vartheta^2>0\)。若估计量 \(T(x)\) 在此区间无偏，即 \(\mathbb E_\vartheta T=\vartheta\)，且方差有限，则 \(\operatorname{Var}_\vartheta T\geq I(\vartheta)^{-1}\)。
\end{proposition}
\begin{proof}
由概率归一化求导，\(\mathbb E S_\vartheta=\sum_x\partial_\vartheta P_\vartheta(x)=0\)。由无偏条件求导并用 \(\partial_\vartheta P=P S_\vartheta\)，有 \(1=\mathbb E[(T-\vartheta)S_\vartheta]\)。对随机变量 \(T-\vartheta\) 与 \(S_\vartheta\) 用 Cauchy--Schwarz 不等式，得到 \(1\leq\sqrt{\operatorname{Var}T}\sqrt{I(\vartheta)}\)，平方即证。若 \(P_\vartheta(x)=0\) 或参数处在边界，需要另行检查这些正则条件，不能直接套用此式。
\end{proof}

若同时估计多个参数，把各偏导排成得分列向量 \(S\)，信息矩阵就是 \(I=\mathbb E(SS^\top)\)。在同样的正则条件下，对无偏估计列向量 \(T\) 求导可得 \(\mathbb E[(T-\vartheta)S^\top]=\mathbf1\)。若 \(I\) 可逆，随机向量 \(Z=T-\vartheta-I^{-1}S\) 的协方差矩阵半正定。展开并用刚才的恒等式，便得到矩阵界 \(\operatorname{Cov}(T)-I^{-1}=\operatorname{Cov}(Z)\succeq0\)。这里 \(A\succeq0\) 的定义是对每个实向量 \(u\) 都有 \(u^\top Au\geq0\)；这也解释了为何信息矩阵不可逆时不能对不可识别方向声称有限误差界。

对只估计相位、把 \(r>0\) 当作已知的局部实验，令 \(\vartheta=\phi\)，刚才的信息给出 \(\operatorname{Var}\widehat\phi\gtrsim1/(\nu r^2)\)。若 \(r\) 也未知，则用上述二维矩阵界；由于 \(r\) 与 \(\phi\) 的导数列在小 \(r\) 时正交，其领先阶仍是 \(I_{\phi\phi}=\nu r^2\)。此界只约束满足所述无偏与正则条件的估计量，绝非保证某个算法必能达到它。

一般地，参数 \(\vartheta=(\vartheta^1,\ldots,\vartheta^p)\) 到结果概率 \(p_j(\vartheta,s)\) 的映射，若不同参数始终产生不同的所有设置概率，称全局可识别；其 Jacobian \(J_{(j,s),a}=\partial p_j(\vartheta,s)/\partial\vartheta^a\) 满列秩是局部一阶可识别的充分条件。秩不足时存在参数微扰 \(u\neq0\) 使 \(Ju=0\)，数据对这一方向的一阶变化为零。

若设置 \(s\) 的预期入射电子数为 \(\nu_s\)，探测通道 \(j\) 的平均计数为 \(\mu_{js}(\vartheta)=\nu_s p_j(\vartheta,s)\)。独立 Poisson 计数 \(N_{js}\) 的对数似然为
\begin{equation}
\log\mathcal L(\vartheta)=\sum_{j,s}
\bigl[N_{js}\log\mu_{js}(\vartheta)-\mu_{js}(\vartheta)-\log(N_{js}!)\bigr].
\label{eq:feqo-count-likelihood}
\end{equation}
由 \(\partial_a\log\mathcal L=\sum (N/\mu-1)\partial_a\mu\)、\(\mathbb E N=\mu\)、\(\operatorname{Var}N=\mu\)，得到 Fisher 信息
\begin{equation}
I_{ab}(\vartheta)=\sum_{j,s}\frac{\partial_a\mu_{js}(\vartheta)
\,\partial_b\mu_{js}(\vartheta)}{\mu_{js}(\vartheta)}.
\label{eq:feqo-poisson-fisher}
\end{equation}
若 \(Ju=0\)，同时有 \(u^\top Iu=0\)。增加曝光时间只放大已有非零信息，不能填补严格零方向。若总计数被实验固定，应改用多项分布似然，不能把相关的分箱计数又当独立 Poisson。

即使 Jacobian 满秩，极小奇异值也会使噪声被放大。对线性化模型 \(y=Ax+\epsilon\)，令 \(W\) 为噪声权重，\(Lx\) 为希望平滑或压小的量。Tikhonov 估计最小化 \(\|Ax-y\|_W^2+\lambda\|Lx\|^2\)，其一阶条件为
\begin{equation}
(A^\dagger WA+\lambda L^\dagger L)\widehat x_\lambda=A^\dagger Wy.
\label{eq:feqo-tikhonov}
\end{equation}
在无噪声 \(y=Ax_*\) 下，只要左矩阵可逆，偏差为 \(\widehat x_\lambda-x_*=-\lambda(A^\dagger WA+\lambda L^\dagger L)^{-1}L^\dagger Lx_*\)。正则化稳定小奇异值，也引入可计算偏差；\(\lambda\) 应由独立验证数据、预测风险或明确的先验选择，不能称作物理退相干率。

计数涨落可以通过式~\eqref{eq:feqo-count-likelihood} 的重复抽样或信息矩阵估计；束流强度、响应核标定和参考相位的不确定度要另行传播。优化器收敛阈值只衡量数值误差，不是这些实验误差的替身。

\begin{exercise}
对式~\eqref{eq:feqo-reference-phase-probability} 求 \(\theta=0,\pi/2\) 时关于 \((r,\phi)\) 的 Jacobian 行列式，说明它在 \(r=0\) 为何为零。
\end{exercise}

